%-------------------------
% Resume in LaTeX
% Author : Sofian Zalouk
% Based off of: https://github.com/sb2nov/resume
% License : MIT
%-------------------------

\documentclass[letterpaper,11pt]{article}

\usepackage{latexsym}
\usepackage[empty]{fullpage}
\usepackage{titlesec}
\usepackage{marvosym}
\usepackage[usenames,dvipsnames]{xcolor}
\usepackage{verbatim}
\usepackage{enumitem}
\usepackage[hidelinks]{hyperref}
\usepackage{fancyhdr}
\usepackage[english]{babel}
\usepackage{tabularx}
\ifdefined\pdfgentounicode\input{glyphtounicode}\fi

\hypersetup{
  pdftitle={Zalouk, Sofian Resume},
  pdfauthor={Sofian Zalouk}
}

% Red placeholder for anything still to be filled in. Search for \fillin.
\newcommand{\fillin}[1]{\textcolor{red}{\textbf{[#1]}}}

\pagestyle{fancy}
\fancyhf{}
\fancyfoot{}
\renewcommand{\headrulewidth}{0pt}
\renewcommand{\footrulewidth}{0pt}

% Adjust margins
\addtolength{\oddsidemargin}{-0.5in}
\addtolength{\evensidemargin}{-0.5in}
\addtolength{\textwidth}{1in}
\addtolength{\topmargin}{-.5in}
\addtolength{\textheight}{1.0in}

\urlstyle{same}
\raggedbottom
\raggedright
\setlength{\tabcolsep}{0in}

% Sections formatting
\titleformat{\section}{
  \vspace{-4pt}\scshape\raggedright\large
}{}{0em}{}[\color{black}\titlerule \vspace{-2pt}]
\titlespacing*{\section}{0pt}{14pt}{6pt}

% Ensure that generated pdf is machine readable/ATS parsable
\ifdefined\pdfgentounicode\pdfgentounicode=1\fi

%-------------------------
% Custom commands
\newcommand{\resumeItem}[1]{
  \item\small{
    {#1 \vspace{-2pt}}
  }
}

\newcommand{\resumeSubheading}[4]{
  \vspace{-2pt}\item
    \begin{tabular*}{0.97\textwidth}[t]{l@{\extracolsep{\fill}}r}
      \textbf{#1} & #2 \\
      \textit{\small#3} & \textit{\small #4} \\
    \end{tabular*}\vspace{-7pt}
}

\newcommand{\resumeSubheadingEdu}[6]{
  \vspace{-2pt}\item
    \begin{tabular*}{0.97\textwidth}[t]{l@{\extracolsep{\fill}}r}
      \textbf{#1} & #2 \\
      \small#3 & \textit{\small #4} \\
      \small#5 &  \\
      \small#6 &  \\
    \end{tabular*}\vspace{-7pt}
}

\newcommand{\resumeSubheadingEduGPA}[5]{
  \vspace{-2pt}\item
    \begin{tabular*}{0.97\textwidth}[t]{l@{\extracolsep{\fill}}r}
      \textbf{#1} & #2 \\
      \small#3 & \textit{\small #4} \\
      \small#5 &  \\
    \end{tabular*}\vspace{-7pt}
}

% Publication: {title (with link)}{year}{authors}{venue}
% Uses tabularx so long titles and author lists wrap on their own.
\newcommand{\resumePub}[4]{
  \vspace{-2pt}\item
    \begin{tabularx}{0.97\textwidth}[t]{@{}>{\raggedright\arraybackslash\hyphenpenalty=10000}X@{\hspace{1em}}r@{}}
      \textbf{#1} & #2 \\
      \small #3 & \\
      \small\textit{#4} & \\
    \end{tabularx}\vspace{3pt}
}

\newcommand{\resumeProjectHeading}[2]{
    \item
    \begin{tabular*}{0.97\textwidth}{l@{\extracolsep{\fill}}r}
      \small#1 & #2 \\
    \end{tabular*}\vspace{-7pt}
}

\renewcommand\labelitemii{$\vcenter{\hbox{\tiny$\bullet$}}$}

\newcommand{\resumeSubHeadingListStart}{\begin{itemize}[leftmargin=0.15in, label={}]}
\newcommand{\resumeSubHeadingListEnd}{\end{itemize}}
\newcommand{\resumeItemListStart}{\begin{itemize}[itemsep=1pt, topsep=3pt]}
\newcommand{\resumeItemListEnd}{\end{itemize}\vspace{0pt}}

%-------------------------------------------
%%%%%%  RESUME STARTS HERE  %%%%%%%%%%%%%%%%%
\begin{document}

%----------HEADING----------
\begin{center}
    \textbf{\Huge \scshape Sofian Zalouk} \\ \vspace{3pt}
    \small
    \href{mailto:saz43@cornell.edu}{\underline{saz43@cornell.edu}} $|$
    \href{https://szalouk.github.io/}{\underline{szalouk.github.io}} $|$
    \href{https://scholar.google.com/citations?hl=en&user=dmQmmUkAAAAJ}{\underline{Google Scholar}}
\end{center}


%-----------RESEARCH INTERESTS-----------
\section{Research Interests}
\small{Generative modeling of language, especially diffusion language models: architectures, training objectives, and sampling methods that let models revise their own outputs and learn efficiently from data. Broader interests in probabilistic modeling and uncertainty quantification.}

%-----------EDUCATION-----------
\section{Education}
  \resumeSubHeadingListStart
    \resumeSubheadingEduGPA
      {Cornell University}{2024 -- Present}
      {Ph.D. in Computer Science}{Ithaca, NY}
      {Advisors: Christopher De Sa, Kilian Q. Weinberger}
      \vspace{2pt}

    \resumeSubheadingEdu
      {Stanford University}{2024}
      {M.S. in Computer Science, Concentration: Artificial Intelligence}{Stanford, CA}{Advisors: Stefano Ermon, Andrew Ng}{GPA: 4.14/4.00}
      \vspace{1pt}

    \resumeSubheadingEdu
      {University of Toronto}{2020}
      {B.A.Sc. in Electrical Engineering}{Toronto, Canada}{Ranked 1\textsuperscript{st} out of over 100 students in department.}{GPA: 3.96/4.00}
  \resumeSubHeadingListEnd


%-----------PUBLICATIONS-----------
\section{Publications}
  \resumeSubHeadingListStart

    \resumePub
      {\href{https://arxiv.org/abs/2611.19127}{Empowering Masked Diffusion Models to Self-Correct with Leave-One-Out Transformers}}{2026}
      {\textbf{Sofian Zalouk}, Vincent Counathe, Paul J\"unger, Daniel Cao, Daniel Zlotnick, Kilian Q. Weinberger, Christopher De Sa}
      {NeurIPS -- Conference on Neural Information Processing Systems, 2026}

    \resumePub
      {\href{https://arxiv.org/abs/2505.15962}{Pre-training Limited Memory Language Models with Internal and External Knowledge}}{2026}
      {Linxi Zhao, \textbf{Sofian Zalouk}, Christian K. Belardi, Justin Lovelace, Jin Peng Zhou, Ryan Thomas Noonan, Dongyoung Go, Kilian Q. Weinberger, Yoav Artzi, Jennifer J. Sun}
      {ICLR -- International Conference on Learning Representations, 2026 \textup{\textbf{(Top 2\%)}}\newline
       \textup{\textbf{Best Paper Runner-up}}, CCFM Workshop @ NeurIPS 2025}

    \resumePub
      {\href{https://arxiv.org/abs/2602.20528}{Stop-Think-AutoRegress: Language Modeling with Latent Diffusion Planning}}{2025}
      {Justin Lovelace, Christian Belardi, \textbf{Sofian Zalouk}, Adhitya Polavaram, Srivatsa Kundurthy, Kilian Q. Weinberger}
      {COLM -- Conference on Language Modeling, 2025}

    \resumePub
      {\href{https://openreview.net/forum?id=kBEwAecyLa}{Multivariate Calibration is Performative: A Perspective on Pitfalls and Progress}}{2025}
      {\textbf{Sofian Zalouk}, Charles Marx, Syrine Belakaria, Christopher De Sa, Stefano Ermon}
      {ICML Workshop on Foundation Models for Structured Data, 2025}

    \resumePub
      {\href{https://arxiv.org/abs/2310.20211}{Calibration by Distribution Matching: Trainable Kernel Calibration Metrics}}{2023}
      {Charles Marx$^*$, \textbf{Sofian Zalouk$^*$}, Stefano Ermon \quad ($^*$Equal contribution)}
      {NeurIPS -- Conference on Neural Information Processing Systems, 2023}

    \resumePub
      {\href{https://ieeexplore.ieee.org/document/10337735}{A System for Automated Vehicle Damage Localization and Severity Estimation Using Deep Learning}}{2023}
      {Yuntao Ma, Hiva Ghanbari, Tianyuan Huang, Jeremy Irvin, Oliver Brady, \textbf{Sofian Zalouk}, Hao Sheng, Andrew Ng, Ram Rajagopal, Mayur Narsude}
      {IEEE Transactions on Intelligent Transportation Systems, 2023}

  \resumeSubHeadingListEnd


%-----------RESEARCH EXPERIENCE-----------
\section{Research Experience}
  \resumeSubHeadingListStart

    \resumeSubheading
      {Graduate Research Assistant}{2024 -- Present}
      {Cornell University, with Kilian Q. Weinberger and Christopher De Sa}{Ithaca, NY}
      \resumeItemListStart
        \resumeItem{\textbf{Self-correcting diffusion language models.} Masked diffusion models cannot revise a token once it is generated, and learn from only a fraction of positions. Developing architectures and training objectives that supervise every token and let the model revise its own outputs (LOOT, first author)}
        \resumeItem{\textbf{Hierarchical discrete diffusion.} Organized the vocabulary into a semantic tree rooted at the mask token, so generation refines meaning from coarse to specific and self-correction means moving to a different branch; derived a simplified objective for efficient training}
        \resumeItem{\textbf{Planning in continuous latent space.} Diffusion over sentence embeddings as a planning step for autoregressive generation, with plans steered by lightweight classifiers (STAR-LDM)}
        \resumeItem{\textbf{Separating knowledge from language ability.} Pre-training language models that store facts in an external database, making knowledge editable and verifiable; built the training pipeline and unlearning evaluations (LMLM)}
        \resumeItem{\textbf{Calibrated multivariate forecasts.} Showed that common definitions of multivariate calibration are performative: recalibrating a forecaster changes the test used to judge it (first author, ICML workshop)}
      \resumeItemListEnd

    \resumeSubheading
      {Graduate Research Assistant}{Sep 2020 -- 2024}
      {Stanford Artificial Intelligence Laboratory}{Stanford, CA}
      \resumeItemListStart
        \resumeItem{Uncertainty quantification and calibration with Stefano Ermon (NeurIPS 2023)}
        \resumeItem{Machine learning for healthcare with Andrew Ng and Sharon Zhou}
      \resumeItemListEnd

    \resumeSubheading
      {Research Assistant}{Summer 2022}
      {Ford-Stanford University Collaboration}{Stanford, CA}
      \resumeItemListStart
        \resumeItem{Active learning and image segmentation with Andrew Ng and Ram Rajagopal (IEEE T-ITS 2023)}
      \resumeItemListEnd

    \resumeSubheading
      {Undergraduate Researcher}{Sep 2019 -- May 2020}
      {University of Toronto}{Toronto, Canada}
      \resumeItemListStart
        \resumeItem{Monocular depth estimation with Andreas Moshovos}
      \resumeItemListEnd

  \resumeSubHeadingListEnd


%-----------INDUSTRY EXPERIENCE-----------
\section{Industry Experience}
  \resumeSubHeadingListStart

    \resumeSubheading
      {Software Engineer (Intern)}{Sep 2018 -- Sep 2019}{Intel}{Toronto, Canada}
      \resumeItemListStart
        \resumeItem{Worked on the LLVM compiler to analyze and accelerate deep learning workloads on FPGAs}
        \resumeItem{Improved Intel's OpenCL memory management, reducing runtime and memory overhead by an order of complexity}
      \resumeItemListEnd

    \resumeSubheading
      {ASIC Design Engineer (Intern)}{Summer 2018}{NVIDIA}{Santa Clara, CA}
      \resumeItemListStart
        \resumeItem{Design and verification of processors for deep learning}
      \resumeItemListEnd

  \resumeSubHeadingListEnd


%-----------TEACHING EXPERIENCE-----------
\section{Teaching Experience}
  \resumeSubHeadingListStart
  \vspace{-2pt}\item
     \begin{tabular*}{0.97\textwidth}[t]{l@{\extracolsep{\fill}}r}
      \small Cornell CS 4787/5777 (Principles of Large-Scale Machine Learning), Teaching Assistant; delivered lectures & \small Fall 2026 \\
      \small Stanford CS 228 (Probabilistic Graphical Models), Teaching Assistant  & \small Winter 2024 \\
      \small Stanford CS 229 (Machine Learning), Teaching Assistant  & \small Spring 2021 \\
      \small Stanford CS 236G (Generative Adversarial Networks), Teaching Assistant & \small Winter 2021 \\
    \end{tabular*}\vspace{0pt}
  \resumeSubHeadingListEnd


%-----------AWARDS-----------
\section{Awards and Scholarships}
  \resumeSubHeadingListStart
  \vspace{-2pt}\item
     \begin{tabular*}{0.97\textwidth}[t]{l@{\extracolsep{\fill}}r}
      \small\textbf{Best Paper Runner-up}, CCFM Workshop @ NeurIPS 2025, for Limited Memory Language Models & \small 2025 \\
      \small\textbf{Best Project Award}, Stanford CS 224R (Deep Reinforcement Learning), Instructor: Chelsea Finn & \small Spring 2023 \\
      \small\textbf{W.S. Wilson Medal}, University of Toronto: highest academic standing in the graduating year & \small 2020 \\
      \small\textbf{Outstanding Student Award}, University of Toronto: top 3 students in Electrical Engineering & \small 2015 -- 2020 \\
      \small\textbf{Baptie Scholarship}, University of Toronto: high academic standing (top 1\%) & \small 2016 \\
    \end{tabular*}\vspace{0pt}
  \resumeSubHeadingListEnd


%-----------SKILLS-----------
\section{Technical Skills}
 \begin{itemize}[leftmargin=0.15in, label={}]
    \small{\item{
     \textbf{Languages}{: Python, C/C++, CUDA, JavaScript, HTML/CSS, R, Bash, MATLAB} \\
     \textbf{ML}{: PyTorch, TensorFlow, Keras, scikit-learn, pandas, NumPy, Matplotlib} \\
     \textbf{Tools}{: \LaTeX, Git, Jupyter, GDB, Valgrind}
    }}
 \end{itemize}

\end{document}